%%=============================================================================
%% Conclusie
%%=============================================================================

\chapter{Conclusie}%
\label{ch:conclusie}

% TODO: Trek een duidelijke conclusie, in de vorm van een antwoord op de
% onderzoeksvra(a)g(en). Wat was jouw bijdrage aan het onderzoeksdomein en
% hoe biedt dit meerwaarde aan het vakgebied/doelgroep? 
% Reflecteer kritisch over het resultaat. In Engelse teksten wordt deze sectie
% ``Discussion'' genoemd. Had je deze uitkomst verwacht? Zijn er zaken die nog
% niet duidelijk zijn?
% Heeft het onderzoek geleid tot nieuwe vragen die uitnodigen tot verder 
%onderzoek?

In deze bachelorproef werd gevraagd hoe een veilige, representatieve en reproduceerbare labo-omgeving ontworpen kon worden om studenten Toegepaste Informatica te trainen in de detectie en analyse van traditionele (LKM) en moderne (eBPF) Linux-kernel-rootkits. Deze vraag is beantwoord in de vorm van meerdere .ova-bestanden waarin een klassieke LKM-rootkit (Singularity) en een moderne eBPF-rootkit (bad-bpf) zitten. Om dit te bereiken moest naar de theorie van de beide technieken gekeken worden. LKM-rootkits hebben onbeperkte rechten en maken gebruik van methodes zoals ftrace en directe het geheugen overschrijven. eBPF-rootkits gebruiken compleet andere methodes, het opereert vanuit een veilige virtuele machine binnen de kernel. Het maakt gebruik van Kprobes, Uprobes, en specifieke helper-functies. Daarnaast is onderzocht welke ingebouwde beveiligingen (zoals SMEP, SMAP, KASLR en Secure Boot) een rol spelen bij het installeren van deze rootkits. Tot slot is beschreven hoe systemen kunnen worden beschermd en hoe rootkits kunnen worden opgespoord met behulp van tools zoals Volatility, Falco en Tracee.

Dit onderzoek zorgt ervoor dat het lastige onderwerp van rootkits nu een duidelijke labo-oefening is geworden voor het vak Cybersecurity Advanced. Eerst was het erg moeilijk om rootkits in de les te tonen, omdat ze heel diep in het systeem verborgen zitten. Dankzij de .ova-bestanden en stappenplannen kunnen studenten mee oefenen. Zo zien ze hoe een systeem overgenomen wordt, en leren ze direct hoe ze dit kunnen herstellen en voorkomen.

\section{Kritische reflectie}
Tijdens het opzetten van de labo-omgeving zijn er een aantal uitdagingen naar voren gekomen. Het grootste struikelblok was de gevoeligheid van eBPF-rootkits voor specifieke kernelversies. Hoewel eBPF in theorie gebruikmaakt van CO-RE (Compile Once, Run Everywhere), bleek in de praktijk dat kleine verschillen in de kernel de werking kunnen verstoren. Om deze reden moesten we specifieke kernelpakketten bevriezen (`apt-mark hold`), wat in een echte productieomgeving minder wenselijk is omdat het beveiligingsupdates blokkeert.

Daarnaast was het genereren van de ISF-symbolen voor Volatility 3 via `dwarf2json` erg tijdsintensief en vroeg het veel geheugen van de virtuele machine. Dit is daarom vooraf uitgevoerd voor de studenten. Tot slot bleek dat de eBPF-rootkit actief probeerde te voorkomen dat forensische kernelmodules zoals LiME werden ingeladen. Dit dwong ons om een RAM-dump direct via de hypervisor (VirtualBox) te maken. Hoewel dit uitstekend werkt en het concept van de "onbetrouwbare machine" goed demonstreert, betekent dit wel dat de detectiemethode afhankelijk is van VirtualBox-commando's op de host in plaats van tools binnen de VM zelf.

\section{Toekomstig onderzoek}
Voor toekomstig onderzoek zijn er verschillende interessante richtingen. Ten eerste zou er gekeken kunnen worden naar de werking van eBPF-rootkits op recentere kernelversies (zoals Linux kernel 6.x). De beveiligingsmechanismen van Linux blijven zich ontwikkelen, en het is relevant om te onderzoeken of de besproken rootkits daar nog steeds effectief zijn.

Ten tweede kan er dieper worden ingegaan op de automatische detectie en mitigatie. Het labo richt zich nu vooral op offline geheugenanalyse met Volatility 3. Toekomstig onderzoek kan zich richten op het realtime detecteren of blokkeren van eBPF-malware met behulp van eBPF-beveiligingstools zoals Falco of Tracee, en hoe deze tools zelf kunnen worden beschermd tegen sabotage door een rootkit. Tot slot kan de labo-omgeving in de toekomst geautomatiseerd worden met behulp van Ansible of Vagrant, zodat studenten of docenten de omgeving met één commando kunnen opzetten.